%%%%%%%%%%%%%%%%%%%%%%%%%%%%%%%%%%%%%%%%%%%%%%%%%%%%%%%%%%%%%%%%%%%%%%
% Lớp 10 | Bất phương trình bậc 2 một ẩn | Giải bất phương trình bậc 2 một ẩn
% Dạng 0D7?2-4: Hệ bất phương trình BPT bậc 2
%%%%%%%%%%%%%%%%%%%%%%%%%%%%%%%%%%%%%%%%%%%%%%%%%%%%%%%%%%%%%%%%%%%%%%


% nguồn: TuLuan_Lop10_Dot2 (nhập 14/11/2026)
\begin{bt}%[0D7V2-4]%[Tex hoa sach Toan 10]%[Đỗ Đường Hiếu]%[ma:95946d16]
	Tìm $D=E\cap G$ biết $E$ và $G$ lần lượt là tập nghiệm của hai bất phương trình trong mỗi trường hợp sau
	\begin{multicols}{2}
		\begin{enumerate}
			\item $2x+3\ge 0$ và $-x+5\ge 0$;
			\item $x+2>0$ và $2x-9<0$.
		\end{enumerate}
	\end{multicols}
	\loigiai{
		\begin{enumerate}
			\item Ta có
			\begin{itemize}
				\item $2x+3\ge 0 \Leftrightarrow x\ge -\dfrac{3}{2}$, suy ra $E=\left[-\dfrac{3}{2}; +\infty \right)$.
				\item $-x+5\ge 0 \Leftrightarrow x\le 5$, suy ra $G=\left(-\infty; 5\right]$.
			\end{itemize}
			Do đó $D=E\cap G=\left[-\dfrac{3}{2}; 5 \right]$.
			\item Ta có
			\begin{itemize}
				\item $x+2>0 \Leftrightarrow x>-2$, suy ra $E=\left(-2; +\infty \right)$.
				\item $2x-9<0 \Leftrightarrow x<\dfrac{9}{2}$, suy ra $G=\left(-\infty; \dfrac{9}{2}\right)$.
			\end{itemize}
			Do đó $D=E\cap G=\left(-2; \dfrac{9}{2}\right)$.
		\end{enumerate}
	}	
\end{bt}
